\label{chap:discussion}
Chapter~\ref{chap:evaluation} reported participant behaviour and feedback regarding the evaluated artefact. This chapter abstracts from these individual sessions to evaluate the underlying design principles and connect the findings to existing literature. Specifically, it states what the evaluation says about each of the design's assumptions and what follows from it, articulates the implications for hybrid planning interfaces, and situates the results within the prior work reviewed in Chapter~\ref{chap:related-work}. It then delineates the methodological limitations bounding the findings and asks how far the artefact and the findings carry beyond this study. The formal answers to the research questions and directions for future work are reserved for Chapter~\ref{chap:conclusion}. 
Following the nested model's premise that downstream validation can expose upstream flaws, it also reads those findings back onto the levels above, including the level~1 assumption that a combined graph-and-table interface is the right approach (Chapter~\ref{chap:methodology}, Section~\ref{sec:paradigm}).
Ultimately, the analysis presented here is exploratory. Rather than declaring final laws of user behaviour, it synthesises the observed sessions into hypotheses, offering a foundation for larger-scale studies to test and refine.

\section{What the Evaluation Says About the Design}
\label{sec:disc-assumptions}
\label{sec:disc-interpretation}

The design of Chapter~\ref{chap:design} rested on assumptions it could not test, some inherited from the formative study and some made as design judgements. Each is taken in turn: what was assumed, what the evaluation says about it, and what follows. The loop itself comes first, since each assumption is read against it. The order is the planning loop's, outwards from the surface that drove it. Where a finding yields a rule for other tools, the rule is stated in Section~\ref{sec:disc-implications}.

\paragraph{The result is a process, not a feature.}
Participants converged on one stable loop, and that loop, rather than any single component, is what their sessions were organised around: the pre-fill seeds a plan, the checklist drives it, the recommendations, the catalogue and the graph propose candidates, the Parking Stage holds them, and the table commits them. It is the same for the most fluent participants and for the most junior, who had to be shown where it started, and it is measurable, its two-station cycle accounting for 44.9\,\% of observed state changes in Scenario~A and 38.6\,\% in Scenario~B. What it externalises is the part of programme planning students cannot hold in their heads, which requirements remain open and what would close them, so that attention goes to choosing rather than to tracking.

\paragraph{The Dashboard as the student's compass.}
Held, and it did more than intended. The missing-requirements list drove the process for every participant (E-G15, 11/11), and because the briefs stated no curriculum rules, participants met the curriculum only through it (Section~\ref{sec:eval-conditions}): for the duration of a task the tool did not support an understanding of the curriculum, it constituted it. What it did not hold them to therefore did not exist for them. Terminology it left undefined stayed unclear (E-P01); the two least confident participants used it and the recommendations in place of curriculum knowledge (E-N06); and the reduced-load semester, the one constraint of both briefs it cannot model, is the unmet constraint in half of the runs that missed their brief (Section~\ref{sec:eval-effectiveness}).

\paragraph{A staging area between finding and committing.}
Held, and more important than the design claimed. Every participant shortlisted before committing (E-G06, 11/11), and the Parking Stage is what lets two views work one plan: a buffer outside the plan's accounting where courses found in the graph, or surfaced by a recommendation, wait until the table assigns them. A shared-plan hybrid does not need each view to do everything, but it does need an explicit place for a decision made and not yet committed. Its cost was not anticipated: parked credits are invisible to every completion readout (E-P38, 10/11), and several participants compensated for that in their heads.

\paragraph{Graph for exploration, table for commitment.}
Supported, but narrower than its wording and true for a different reason than the design supposed. The graph served the query step and the table the commit step for every participant who used both (Section~\ref{sec:eval-graphrole}), though what students did in the graph was directed querying, turning a named gap into a candidate set, rather than the open-ended browsing the formative study called exploration (Section~\ref{sec:mutual-exclusion}). It does not add a station to the cycle; it takes over the one the catalogue held, at the catalogue's occupancy and entry count, leaving the cycle's shape and period unchanged (Section~\ref{sec:eval-substitution}). Nor is the division imposed by the representations themselves: plan changes were made inside the graph by participants across the experience range, and what sends placement back to the table is that semester lanes and their totals are not rendered in Graph View, so a student placing from a node cannot see the load consequence, and one breached a reduced-load cap that way. Within the query step the work was done largely by the filter sidebar rather than the node--link drawing: participants said so (E-G07, 11/11; E-P04, 11/11), including those who read the graph most fluently, and one entered the graph for the single purpose of finding a course of a given ECTS value, a query the catalogue cannot express. Interest-driven browsing (E-G10, 9/11) and module composition (E-G17, 6/11) are things the filters do not give, but the study cannot tell how much of the graph's contribution would survive if the table had the filters (Chapter~\ref{chap:conclusion}, Section~\ref{sec:concl-future}). The claim is about a structural view entering an established workflow rather than one met from the start; Section~\ref{sec:disc-related} reads the seam through cognitive fit.

\paragraph{An overlay that draws only what the curriculum fixes.}
Sound as a rule, and empty in application. An edge should carry a consequence, which is what keeps an overlay legible, but the Bachelor programme the study planned in fixes no ordering between two named courses at all, the two it publishes being advisory, which the rule excluded. The overlay in front of participants was therefore empty rather than merely sparse, and four of them asked for ordering information at the point of placement (E-P44, 4/11). Section~\ref{sec:disc-related} draws the consequence for the graph-visualisation literature.

\paragraph{Guide rather than block.}
Held for the rules it covered (E-G04, 11/11), with two problems in which rules those were and in how tightly they bound. The workload band is split across the two tiers: exceeding the ceiling is blocked, while falling below the floor raises only a dismissible warning, so an under-loaded semester is never forced to correct (E-P35, 10/11). And the band's thresholds come from the student's profile rather than from the curriculum, so where the Dashboard's limit was a rule it cannot model, this is one it models against the wrong numbers: two Scenario~B runs ended with a semester above the brief's ceiling and no warning, because the ceiling in force was still the default (Section~\ref{sec:eval-config}). The setting is discoverable and it works; the consequence is structural rather than about particular students. A checker whose thresholds come from the user is correct only for those who supply them and reports the same confident verdict either way, so a plan checked against the wrong band is indistinguishable from one checked against the right one. Which students are at risk is not something eleven sessions can establish: both unconfigured runs are by participants who reported no prior experience with study-planning tools, but four other participants who reported none did set the profile.

\paragraph{Recommendations with reasons, beside the plan.}
Held for use and half held for reasons. Reasoned suggestions were read and used by ten of eleven participants, but they did not drive the process, and whom they served divided by curriculum confidence: those who could not yet judge the curriculum leaned on them in place of domain knowledge, while those who could browsed and returned to the checklist and catalogue search (E-N05, 7/11), although the engine visibly responded to their stated interests. An explainable recommendation is therefore worth most to the least expert user, and confident users want a different route to the same end. The reasons were thinner than intended as well: a card named the relationship behind a suggestion, that a course resembles one already in the plan or that similar students chose it, but not which of the student's own profile entries had produced it, so the evidence could be read and not traced (E-P68, 4 of the 9 participants for whom the code could be decided).

\paragraph{Self-report measured felt completion.}
Eight of the 21 resolved runs missed their brief; seven of those eight were rated a success by the participant who produced them, and all twenty-two runs ended with the Dashboard reporting no missing requirements, a report correct for the curriculum's requirements and silent on the task's (Section~\ref{sec:eval-effectiveness}). A student can hold a correct model of the brief in their head and lose it, because that model and the checklist are never weighed against each other: the checklist speaks in the interface, continuously and with authority, while the student's own constraint has no representation there and has to be remembered against it. When the two conflict the checklist prevails: it is continuously present and the student's constraint is not, so the constraint that is read is the constraint that is satisfied. A constraint with no representation in the interface cannot compete with one that has it.

\paragraph{The two upstream assumptions.}
Above these sit the two assumptions the design inherited rather than made (Section~\ref{sec:overall-assumptions}). The first is that a combined graph-and-table interface is worth building at all. On this the evaluation says only that the two surfaces served different steps of one loop and that seven participants named that division as the value of having both (E-N01, 7/11); no non-hybrid alternative was run, and Chapter~\ref{chap:conclusion} (Section~\ref{sec:concl-future}) states the open question. The second is that prerequisite sequencing is the planning-relevant structure of a curriculum. Here the finding is the curriculum's rather than the evaluation's: the regulations in scope publish almost no sequencing, and students used the graph for what such a curriculum affords: finding candidates, browsing by interest (E-G10, 9/11) and reading how modules are composed (E-G17, 6/11).


\section{Design Implications}
\label{sec:disc-implications}
The findings above yield six rules for interfaces that guide a user through a constrained, multi-step decision, of which a study planner is one instance. Each rule is stated for that class, followed by the form it took in this tool and the finding it rests on; the argument for each is made in Section~\ref{sec:disc-assumptions} and is not repeated here. The rules are written to hold beyond this artefact.

\begin{itemize}
  \item \textbf{An authoritative signal must cover every actual constraint.} Users do not read a partial completion signal as partial. If an interface reports a plan as complete, it must compute that status against all rules the user is accountable for, not just the ones the system easily models. Pending work should be explicitly counted or excluded, and thresholds should default to domain standards (warning the user if defaults are active). In this tool: the system ignored the brief's reduced-load requirement and invisible parked credits.
  
  \item \textbf{Tier rules by the cost of the breach.} Whether a rule firmly blocks an action or merely warns the user should depend on the actual penalty for getting it wrong. Opposing bounds of the same constraint (e.g., minimum vs.\ maximum load) should enforce the same level of friction. In this tool: the workload ceiling actively blocked users, but the floor only raised a dismissible warning.
  
  \item \textbf{Allow commitment only from views that show the consequences.} If a surface lets a user commit an action, it must display the state that action alters. If a view cannot show those consequences, it should be restricted to finding and staging candidates. In this tool: the graph lacks a semester load dimension, so placing a course from the graph hides the resulting workload.
  
  \item \textbf{Put decision-making data on the surface where the decision happens.} Separating a choice from its evidence forces a round trip, which users only make if they remember to. Controls and status updates must live exactly where the user acts on them. In this tool: filters lived in the graph, but placement happened in the table. Requirement statuses were hidden behind collapsible panels.
  
  \item \textbf{Show valid targets before an action is attempted.} If the system knows which moves are legal, it should highlight them before the user acts, rather than rejecting invalid attempts afterward. In this tool: legal semester lanes were not highlighted during drag-and-drop, and a course's valid term was not visible on the lane.
  
  \item \textbf{Treat guided and self-directed workflows as equals.} Novices rely on guidance while experts bypass it; neither is simply a fallback for the other. Both need equal prominence. Furthermore, automated guidance must explain its inputs so confident users can audit the logic instead of outright discarding it. In this tool: recommendations competed with the checklist and manual search. Explanations stated that a relationship existed, but not which user profile entry caused it.
\end{itemize}

\section{Relation to Prior Work}
\label{sec:disc-related}

The findings are consistent with, qualify and extend the work reviewed in Chapter~\ref{chap:related-work}. The three are taken in turn, and the extension is developed over the five paragraphs that follow it: the dashboard critique, the integration this tool was built to test, and the theory that accounts for what the graph did.

First, the results are \emph{consistent with} the specific usability friction points identified by Bartel et al.~\cite{bartel_design_2024} regarding study-planning assistants. In a different interface and sample, participants in this study met the same two difficulties, reading institutional curriculum terminology and finding nested filter controls. Design principles name known hazards without conferring immunity to them.

Second, the findings \emph{qualify} the foundational premise of the graph-visualisation literature. Systems such as \emph{ViCurriAS}~\cite{zucker_vicurrias_2009}, \emph{STOPS}~\cite{auvinen_stops_2014}, and the elective visualisations by Trippel and Röpke~\cite{trippel_developing_2025} rest on the premise that a curriculum carries a prerequisite network worth drawing. The curricula in this study carry little of one. Outside the introductory phase the regulations state few dependencies, so the implemented graph draws containment rather than sequence, and a prerequisite advisory arose once across the twenty-two runs. For curricula that encode few hard dependencies, the premise that a prerequisite network is what a curriculum graph is for therefore does not hold. What such a graph is worth instead is a question about use rather than about the regulations, and the paragraph that follows puts the evaluation's answer to it.

Third, the results \emph{extend} the AIStudyBuddy lineage~\cite{roepke_study_2024, judel_supporting_2023, trippel_developing_2025} by capturing something this prior work does not: the step-by-step process of a student actually forming a plan. Existing studies in this lineage have established that graph-based visualisations are well received~\cite{trippel_developing_2025}, and they have analysed study paths after the fact~\cite{roepke_study_2024}. However, neither approach observes the planning process as it unfolds. The session recordings from this evaluation qualify the former and reach past the latter. While positive reception is reproduced here, it does not predict the graph's actual role: participants used the graph as a query surface that displaced the catalogue, rather than as a primary planning surface, and those who read it most fluently judged it substitutable by its own filter panel. Furthermore, watching the live process reveals behaviours that post-hoc institutional metrics cannot capture. It shows exactly how students rely on completion signals and recommendations while working, and it exposes the critical gap between a tool reporting a plan as complete and the plan actually meeting all constraints.

This gap gives the dashboard critique of Chapter~\ref{chap:related-work} a counterpart at the level of the individual student. Caulfield's objection to Course Signals~\cite{caulfield2013} was that its retention result was read as an outcome the signal had produced when the analysis had not established that it had: the signal was credited beyond what had been checked. The completion finding here has the same shape at a smaller scale. The Dashboard reported no missing requirements, correctly for the curriculum's rules and silently on the task's, and participants read that report as completion in preference to their own model of the task. Parasuraman and Riley~\cite{parasuraman1997humans} describe the general hazard as misuse of automation, over-reliance that produces failures of monitoring or decision biases, and name the saliency of a system's state indicators among the factors that govern it. A completion checklist is an automation of that kind: continuous, salient and authoritative, and silent about what it has not checked, so a student who trusts it has no cue that anything remains to be monitored. The parallel is not a claim about Course Signals, whose signal and setting differ; it is that a trusted signal invites the same over-reliance whatever it measures.

These patterns also speak to the integration Bodily and Verbert asked for. Of the 93 articles in their review, 9\,\% describe an enhanced visualisation with recommendations and 6\,\% a visualisation with recommendations and data mining. Counted more broadly, as any article pairing a recommendations component with a visualisation or text feedback, 17\,\% qualify, and they call for interfaces that address both \emph{what} to do and \emph{why}~\cite{bodily_review_2017}. This tool provides both components, and the frame record shows them used together but not as one. The Dashboard drove the planning loop for every participant; the Recommendation Panel accompanied it, but entered a fraction as often as the catalogue, and set aside some participants in favour of the checklist and search once they could judge the curriculum for themselves. The combination is used, with a division of labour that follows curriculum confidence; whether it works better than either component alone is not something this design can test.

That division has a counterpart in the recommender literature. Ma et al.~\cite{ma_courseq_2021} evaluated a course-recommendation interface pairing a search-and-filter component with a visual explanation, and found that students with clear learning goals engaged more with search and filtering, while students without clear goals interacted more with the visualisation and answered significantly positively on perceived usefulness, trust and satisfaction. The split reported here runs along the same line, with the reasoned recommendations in the role their visualisation played: participants who could not yet judge the curriculum leaned on them, and those who could returned to the checklist and search. Interfaces, instruments and samples differ, so the two are consistent rather than mutually confirming.

Finally, the displacement of the catalogue by the graph is an instance of cognitive fit~\cite{vessey1991cognitive}. Vessey's theory holds that problem-solving performance depends on the match between representation and task: spatial tasks, those that require perceiving relationships, are served by graphs, while symbolic tasks, those that require extracting or acting on discrete values, are served by tables. The hybrid interface makes that distinction physical.



\section{Limitations}
\label{sec:disc-limitations}
The evaluation took one direction through a tool that admits many, so the limitations below bound this study design rather than the artefact. Most are the price of a decision made on purpose, to trace a small cohort frame by frame rather than survey a larger one. These boundaries dictate how the findings may be interpreted.

\paragraph{The fixed task order limits causal claims about the graph.} 
Because participants always completed the table-only task before the graph-enabled task, the graph's introduction is entangled with practice effects. While the graph clearly displaced the catalogue (Section~\ref{sec:eval-substitution}), the faster placement rates in the second task may simply be the result of practice. Furthermore, because participants had already established a tabular workflow before the graph was introduced, the observation that the graph acted only as a brief "excursion" from the main loop might be a result of the study's sequence rather than natural user behaviour (Chapter~\ref{chap:evaluation}, Section~\ref{sec:eval-order}).

\paragraph{The findings test this specific implementation, not all hybrid planners.} 
The evaluated tool is just one way to translate the formative study's requirements into a working interface (Section~\ref{chap:from-themes-to-requirements}). Consequently, it is difficult to separate the fundamental limits of hybrid planning from the specific design choices made here.

\paragraph{A small, sample bounds generalisation.} 
The eleven participants were drawn from computing programmes at a single institution. These students are likely more comfortable reading node--link diagrams than the general student body. If a broader population were tested, they would likely find the graph even harder to read, reinforcing the conclusion that the filter panel is the graph's most universally valuable component. Additionally, any patterns based on user expertise rely on very small subsets of this sample, meaning they should be read as indicators rather than definitive proof.

\paragraph{Researcher dual-roles and lab conditions invite bias.} 
Because the author designed the tool, moderated the sessions, and coded the results, the self-reported questionnaire data is highly vulnerable to participants wanting to please the researcher (demand characteristics). This explains why the survey scores sit near the ceiling, and it is why this chapter treats the objective screen recordings as much stronger evidence. Furthermore, participants were planning for a fictional persona under a time limit, and the moderator occasionally had to unblock them. Therefore, the task outcomes are facilitator-assisted, and it remains an open question how these behaviours transfer to independent, unprompted use.

\paragraph{Measurement and coding methods have technical constraints.} 
The video frame analysis, sampled every ten seconds, inherently misses interface interactions shorter than that interval as well as user intend(Section~\ref{sec:eval-framemethod}).

\section{Generalisability}
\label{sec:disc-generalisability}

Two questions of generalisability are separate, and are answered separately: how far the artefact would carry to another institution, and how far the findings would carry to another tool.

\paragraph{The artefact.}
The curriculum abstraction underlying the tool (Chapter~\ref{chap:design}, Section~\ref{sec:curriculum-abstraction}) was read directly off TU~Wien's Bachelor in Computer Science and Master in Software Engineering curriculum regulations rather than designed as an institution-agnostic model. Separating what would transfer by re-parameterisation from what would need redesign gives a checkable split. Programme-agnostic are the four-level hierarchy of programme, exam subject, module and course; the plan model and the two views' division of labour; the checklist-driven loop and the Parking Stage; and the per-semester credit limits, which are the tool's own plan-sanity limits rather than curriculum law. Bound to the two TU~Wien programmes, and encoded as one rule set each (Chapter~\ref{chap:implementation}, Section~\ref{sec:impl-compliance}), are the two obligation taxonomies, the Bachelor's StEOP course sets and its allowance of credits before the gate, its seven focus areas and their alias table, its variant-mixing rule and its transferable-skills cap, and the Master's core sets per exam subject, its subject-modules minimum, its advanced-topics prefixes and its two enforced prerequisite pairs. A third programme therefore means a third rule set, not a re-parameterisation, and how much of the rule sets could be shared is the question Section~\ref{sec:impl-compliance} leaves to that point.

\paragraph{The findings.}
The thesis's separable outputs, stated as contributions in Chapter~\ref{chap:introduction} (Section~\ref{sec:contributions}), generalise differently. The requirements specification and the process findings are informative for study-planning tools in general. The hybrid-specific findings, namely the graph's position in the planning cycle, its dependence on filters, and the Parking Stage's function as a hinge, are evidence for designing hybrid interfaces specifically, not evidence that a hybrid is the right choice for every planner.
